%%%%%%%%%%%%%%%%%%%%%%%%%%%%%%%%%%%%%%%%%%%%%%%%%%%%%%%%%%%%
% 7.8 PROBABILITY TRANSFORMS
% The Composition of Advanced Mathematics
%%%%%%%%%%%%%%%%%%%%%%%%%%%%%%%%%%%%%%%%%%%%%%%%%%%%%%%%%%%%

\section{Probability Transforms}
\label{sec:probability-transforms}

\prerequisite{
Random variables, discrete and continuous distributions, expectation,
moments, functions of random variables, infinite series, power series,
complex numbers, improper integrals, derivatives, and basic transform
methods.
}

\learningobjectives{
By the end of this section, the reader should be able to:
\begin{itemize}
    \item explain why probability transforms are useful;
    \item define probability generating functions;
    \item recover probabilities from probability generating functions;
    \item compute factorial moments from probability generating
          functions;
    \item define moment generating functions;
    \item compute moments from moment generating functions;
    \item use moment generating functions to identify distributions;
    \item derive transforms of sums of independent random variables;
    \item define characteristic functions;
    \item explain why characteristic functions always exist;
    \item compute characteristic functions for standard distributions;
    \item understand uniqueness of characteristic functions;
    \item use characteristic functions for sums and convergence;
    \item define Laplace transforms of nonnegative random variables;
    \item connect Laplace transforms with moment generating functions;
    \item define cumulant generating functions;
    \item compute basic cumulants;
    \item understand additivity of cumulants for independent sums;
    \item distinguish raw moments, factorial moments, and cumulants;
    \item use transforms to derive standard distributional closure
          properties;
    \item recognize transforms as algebraic encodings of probability
          distributions;
    \item understand transform inversion conceptually;
    \item connect probability transforms with generating functions,
          Fourier analysis, asymptotic probability, and stochastic
          processes.
\end{itemize}
}

%%%%%%%%%%%%%%%%%%%%%%%%%%%%%%%%%%%%%%%%%%%%%%%%%%%%%%%%%%%%
\subsection{Why Use Probability Transforms?}
\label{subsec:why-probability-transforms}
%%%%%%%%%%%%%%%%%%%%%%%%%%%%%%%%%%%%%%%%%%%%%%%%%%%%%%%%%%%%

A probability distribution can often be encoded inside another
function.

This may convert difficult probability operations into simpler
algebraic operations.

For example,

\[
\boxed{
\text{convolution of distributions}
}
\]

often becomes

\[
\boxed{
\text{multiplication of transforms}.
}
\]

Similarly, moments can often be recovered from derivatives of a
transform.

Thus probability transforms provide a bridge between

\[
\boxed{
\text{probability}
+
\text{algebra}
+
\text{analysis}.
}
\]

%%%%%%%%%%%%%%%%%%%%%%%%%%%%%%%%%%%%%%%%%%%%%%%%%%%%%%%%%%%%
\subsection{Main Probability Transforms}
\label{subsec:main-probability-transforms}
%%%%%%%%%%%%%%%%%%%%%%%%%%%%%%%%%%%%%%%%%%%%%%%%%%%%%%%%%%%%

Important transforms include:

\[
\boxed{
\text{probability generating function}
}
\]

\[
\boxed{
\text{moment generating function}
}
\]

\[
\boxed{
\text{characteristic function}
}
\]

\[
\boxed{
\text{Laplace transform}
}
\]

and

\[
\boxed{
\text{cumulant generating function}.
}
\]

Each transform is useful in a different setting.

%%%%%%%%%%%%%%%%%%%%%%%%%%%%%%%%%%%%%%%%%%%%%%%%%%%%%%%%%%%%
\subsection{Probability Generating Functions}
\label{subsec:probability-generating-functions}
%%%%%%%%%%%%%%%%%%%%%%%%%%%%%%%%%%%%%%%%%%%%%%%%%%%%%%%%%%%%

Probability generating functions are particularly useful for
nonnegative integer-valued random variables.

\begin{definitionbox}{Probability Generating Function}
Let \(X\) take values in

\[
\{0,1,2,\ldots\}.
\]

The probability generating function, abbreviated PGF, is

\[
\boxed{
G_X(s)
=
\mathbb{E}[s^X].
}
\]
\end{definitionbox}

For a discrete random variable,

\[
\boxed{
G_X(s)
=
\sum_{k=0}^{\infty}
P(X=k)s^k.
}
\]

%%%%%%%%%%%%%%%%%%%%%%%%%%%%%%%%%%%%%%%%%%%%%%%%%%%%%%%%%%%%
\subsection{PGF as an Ordinary Generating Function}
\label{subsec:pgf-ordinary-generating-function}
%%%%%%%%%%%%%%%%%%%%%%%%%%%%%%%%%%%%%%%%%%%%%%%%%%%%%%%%%%%%

If

\[
p_k=P(X=k),
\]

then

\[
G_X(s)
=
p_0+p_1s+p_2s^2+\cdots.
\]

Thus the PGF is precisely the ordinary generating function of the
probability sequence

\[
p_0,p_1,p_2,\ldots.
\]

This connects directly with the generating-function methods developed
in Chapter 6.

%%%%%%%%%%%%%%%%%%%%%%%%%%%%%%%%%%%%%%%%%%%%%%%%%%%%%%%%%%%%
\subsection{Recovering Probabilities from a PGF}
\label{subsec:recover-probabilities-pgf}
%%%%%%%%%%%%%%%%%%%%%%%%%%%%%%%%%%%%%%%%%%%%%%%%%%%%%%%%%%%%

If

\[
G_X(s)
=
\sum_{k=0}^{\infty}
p_ks^k,
\]

then

\[
\boxed{
p_k
=
[s^k]G_X(s),
}
\]

where

\[
[s^k]
\]

means ``the coefficient of \(s^k\).''

Equivalently, if the PGF is sufficiently differentiable,

\[
\boxed{
P(X=k)
=
\frac{
G_X^{(k)}(0)
}{
k!
}.
}
\]

%%%%%%%%%%%%%%%%%%%%%%%%%%%%%%%%%%%%%%%%%%%%%%%%%%%%%%%%%%%%
\subsection{Normalization of a PGF}
\label{subsec:pgf-normalization}
%%%%%%%%%%%%%%%%%%%%%%%%%%%%%%%%%%%%%%%%%%%%%%%%%%%%%%%%%%%%

Since

\[
G_X(1)
=
\sum_{k=0}^{\infty}
P(X=k),
\]

we obtain

\[
\boxed{
G_X(1)=1.
}
\]

Also,

\[
G_X(0)
=
P(X=0).
\]

Thus,

\[
\boxed{
G_X(0)=P(X=0).
}
\]

%%%%%%%%%%%%%%%%%%%%%%%%%%%%%%%%%%%%%%%%%%%%%%%%%%%%%%%%%%%%
\subsection{PGF of a Bernoulli Distribution}
\label{subsec:pgf-bernoulli}
%%%%%%%%%%%%%%%%%%%%%%%%%%%%%%%%%%%%%%%%%%%%%%%%%%%%%%%%%%%%

Let

\[
X\sim\operatorname{Bernoulli}(p).
\]

Then

\[
G_X(s)
=
(1-p)s^0+ps^1.
\]

Therefore,

\[
\boxed{
G_X(s)
=
1-p+ps.
}
\]

%%%%%%%%%%%%%%%%%%%%%%%%%%%%%%%%%%%%%%%%%%%%%%%%%%%%%%%%%%%%
\subsection{PGF of a Binomial Distribution}
\label{subsec:pgf-binomial}
%%%%%%%%%%%%%%%%%%%%%%%%%%%%%%%%%%%%%%%%%%%%%%%%%%%%%%%%%%%%

Let

\[
X\sim\operatorname{Bin}(n,p).
\]

Then

\[
G_X(s)
=
\sum_{k=0}^{n}
\binom nk
p^k(1-p)^{n-k}s^k.
\]

By the binomial theorem,

\[
\boxed{
G_X(s)
=
(1-p+ps)^n.
}
\]

%%%%%%%%%%%%%%%%%%%%%%%%%%%%%%%%%%%%%%%%%%%%%%%%%%%%%%%%%%%%
\subsection{Worked Example: Recovering a Binomial Probability}
\label{subsec:worked-pgf-binomial-coefficient}
%%%%%%%%%%%%%%%%%%%%%%%%%%%%%%%%%%%%%%%%%%%%%%%%%%%%%%%%%%%%

\begin{workedexamplebox}{Extracting a Probability from a PGF}

Suppose

\[
X\sim\operatorname{Bin}(4,p).
\]

Then

\[
G_X(s)
=
(1-p+ps)^4.
\]

The probability

\[
P(X=2)
\]

is the coefficient of

\[
s^2.
\]

Expanding using the binomial theorem gives

\[
\begin{answerbox}
\[
\boxed{
P(X=2)
=
\binom42
p^2(1-p)^2.
}
\]
\end{answerbox}

\end{workedexamplebox}

%%%%%%%%%%%%%%%%%%%%%%%%%%%%%%%%%%%%%%%%%%%%%%%%%%%%%%%%%%%%
\subsection{PGF of a Geometric Distribution}
\label{subsec:pgf-geometric}
%%%%%%%%%%%%%%%%%%%%%%%%%%%%%%%%%%%%%%%%%%%%%%%%%%%%%%%%%%%%

Under the convention

\[
X
=
\text{number of trials until first success},
\]

we have

\[
P(X=k)
=
(1-p)^{k-1}p,
\qquad
k\geq1.
\]

Then

\[
G_X(s)
=
\sum_{k=1}^{\infty}
(1-p)^{k-1}ps^k.
\]

Factor:

\[
=
ps
\sum_{j=0}^{\infty}
[(1-p)s]^j.
\]

Therefore,

\[
\boxed{
G_X(s)
=
\frac{
ps
}{
1-(1-p)s
},
}
\]

for values of \(s\) where the series converges.

%%%%%%%%%%%%%%%%%%%%%%%%%%%%%%%%%%%%%%%%%%%%%%%%%%%%%%%%%%%%
\subsection{PGF of a Poisson Distribution}
\label{subsec:pgf-poisson}
%%%%%%%%%%%%%%%%%%%%%%%%%%%%%%%%%%%%%%%%%%%%%%%%%%%%%%%%%%%%

If

\[
X\sim\operatorname{Poisson}(\lambda),
\]

then

\[
G_X(s)
=
\sum_{k=0}^{\infty}
e^{-\lambda}
\frac{\lambda^k}{k!}
s^k.
\]

Thus,

\[
=
e^{-\lambda}
\sum_{k=0}^{\infty}
\frac{(\lambda s)^k}{k!}.
\]

Therefore,

\[
\boxed{
G_X(s)
=
e^{\lambda(s-1)}.
}
\]

%%%%%%%%%%%%%%%%%%%%%%%%%%%%%%%%%%%%%%%%%%%%%%%%%%%%%%%%%%%%
\subsection{Derivatives of a PGF}
\label{subsec:pgf-derivatives}
%%%%%%%%%%%%%%%%%%%%%%%%%%%%%%%%%%%%%%%%%%%%%%%%%%%%%%%%%%%%

Differentiate:

\[
G_X'(s)
=
\sum_{k=1}^{\infty}
kP(X=k)s^{k-1}.
\]

Setting

\[
s=1
\]

gives

\[
\boxed{
G_X'(1)
=
\mathbb{E}[X].
}
\]

Thus the first derivative produces the mean.

%%%%%%%%%%%%%%%%%%%%%%%%%%%%%%%%%%%%%%%%%%%%%%%%%%%%%%%%%%%%
\subsection{Second PGF Derivative}
\label{subsec:second-pgf-derivative}
%%%%%%%%%%%%%%%%%%%%%%%%%%%%%%%%%%%%%%%%%%%%%%%%%%%%%%%%%%%%

Differentiate again:

\[
G_X''(s)
=
\sum_{k=2}^{\infty}
k(k-1)
P(X=k)
s^{k-2}.
\]

Therefore,

\[
\boxed{
G_X''(1)
=
\mathbb{E}[X(X-1)].
}
\]

This is the second factorial moment.

%%%%%%%%%%%%%%%%%%%%%%%%%%%%%%%%%%%%%%%%%%%%%%%%%%%%%%%%%%%%
\subsection{Variance from a PGF}
\label{subsec:variance-from-pgf}
%%%%%%%%%%%%%%%%%%%%%%%%%%%%%%%%%%%%%%%%%%%%%%%%%%%%%%%%%%%%

Since

\[
X^2
=
X(X-1)+X,
\]

we have

\[
\mathbb{E}[X^2]
=
G_X''(1)+G_X'(1).
\]

Therefore,

\[
\boxed{
\operatorname{Var}(X)
=
G_X''(1)
+
G_X'(1)
-
[G_X'(1)]^2.
}
\]

%%%%%%%%%%%%%%%%%%%%%%%%%%%%%%%%%%%%%%%%%%%%%%%%%%%%%%%%%%%%
\subsection{Worked Example: Poisson Moments from the PGF}
\label{subsec:worked-poisson-pgf-moments}
%%%%%%%%%%%%%%%%%%%%%%%%%%%%%%%%%%%%%%%%%%%%%%%%%%%%%%%%%%%%

\begin{workedexamplebox}{Mean and Variance from a PGF}

For

\[
X\sim\operatorname{Poisson}(\lambda),
\]

\[
G_X(s)
=
e^{\lambda(s-1)}.
\]

Differentiate:

\[
G_X'(s)
=
\lambda
e^{\lambda(s-1)}.
\]

Thus,

\[
G_X'(1)=\lambda.
\]

Also,

\[
G_X''(s)
=
\lambda^2
e^{\lambda(s-1)},
\]

so

\[
G_X''(1)=\lambda^2.
\]

Hence,

\[
\operatorname{Var}(X)
=
\lambda^2+\lambda-\lambda^2.
\]

\begin{answerbox}
\[
\boxed{
\mathbb{E}[X]
=
\lambda,
\qquad
\operatorname{Var}(X)
=
\lambda.
}
\]
\end{answerbox}

\end{workedexamplebox}

%%%%%%%%%%%%%%%%%%%%%%%%%%%%%%%%%%%%%%%%%%%%%%%%%%%%%%%%%%%%
\subsection{PGFs of Independent Sums}
\label{subsec:pgf-independent-sums}
%%%%%%%%%%%%%%%%%%%%%%%%%%%%%%%%%%%%%%%%%%%%%%%%%%%%%%%%%%%%

Suppose \(X\) and \(Y\) are independent nonnegative integer-valued
random variables.

Then

\[
G_{X+Y}(s)
=
\mathbb{E}[s^{X+Y}].
\]

Since

\[
s^{X+Y}
=
s^Xs^Y,
\]

independence gives

\[
\boxed{
G_{X+Y}(s)
=
G_X(s)G_Y(s).
}
\]

Thus convolution becomes multiplication.

%%%%%%%%%%%%%%%%%%%%%%%%%%%%%%%%%%%%%%%%%%%%%%%%%%%%%%%%%%%%
\subsection{Worked Example: Sum of Poisson Variables}
\label{subsec:worked-pgf-poisson-sum}
%%%%%%%%%%%%%%%%%%%%%%%%%%%%%%%%%%%%%%%%%%%%%%%%%%%%%%%%%%%%

\begin{workedexamplebox}{Poisson Closure by PGFs}

Suppose

\[
X\sim\operatorname{Poisson}(\lambda_1)
\]

and

\[
Y\sim\operatorname{Poisson}(\lambda_2)
\]

independently.

Then

\[
G_X(s)
=
e^{\lambda_1(s-1)}
\]

and

\[
G_Y(s)
=
e^{\lambda_2(s-1)}.
\]

Therefore,

\[
G_{X+Y}(s)
=
e^{(\lambda_1+\lambda_2)(s-1)}.
\]

This is the PGF of a Poisson distribution with parameter

\[
\lambda_1+\lambda_2.
\]

\begin{answerbox}
\[
\boxed{
X+Y
\sim
\operatorname{Poisson}(\lambda_1+\lambda_2).
}
\]
\end{answerbox}

\end{workedexamplebox}

%%%%%%%%%%%%%%%%%%%%%%%%%%%%%%%%%%%%%%%%%%%%%%%%%%%%%%%%%%%%
\subsection{Moment Generating Functions}
\label{subsec:moment-generating-functions}
%%%%%%%%%%%%%%%%%%%%%%%%%%%%%%%%%%%%%%%%%%%%%%%%%%%%%%%%%%%%

Moment generating functions are useful when moments and sums are of
primary interest.

\begin{definitionbox}{Moment Generating Function}
The moment generating function, abbreviated MGF, of \(X\) is

\[
\boxed{
M_X(t)
=
\mathbb{E}[e^{tX}],
}
\]

for those values of \(t\) where the expectation exists.
\end{definitionbox}

At \(t=0\),

\[
\boxed{
M_X(0)=1.
}
\]

%%%%%%%%%%%%%%%%%%%%%%%%%%%%%%%%%%%%%%%%%%%%%%%%%%%%%%%%%%%%
\subsection{Why MGFs Generate Moments}
\label{subsec:why-mgfs-generate-moments}
%%%%%%%%%%%%%%%%%%%%%%%%%%%%%%%%%%%%%%%%%%%%%%%%%%%%%%%%%%%%

Expand

\[
e^{tX}
=
1+tX+\frac{t^2X^2}{2!}+\cdots.
\]

Formally taking expectations,

\[
M_X(t)
=
1
+
t\mathbb{E}[X]
+
\frac{t^2}{2!}\mathbb{E}[X^2]
+\cdots.
\]

Thus the coefficients contain the raw moments.

%%%%%%%%%%%%%%%%%%%%%%%%%%%%%%%%%%%%%%%%%%%%%%%%%%%%%%%%%%%%
\subsection{MGF Derivatives}
\label{subsec:mgf-derivatives}
%%%%%%%%%%%%%%%%%%%%%%%%%%%%%%%%%%%%%%%%%%%%%%%%%%%%%%%%%%%%

When differentiation under the expectation is justified,

\[
M_X'(t)
=
\mathbb{E}[Xe^{tX}].
\]

At

\[
t=0,
\]

\[
\boxed{
M_X'(0)
=
\mathbb{E}[X].
}
\]

More generally,

\[
\boxed{
M_X^{(k)}(0)
=
\mathbb{E}[X^k].
}
\]

%%%%%%%%%%%%%%%%%%%%%%%%%%%%%%%%%%%%%%%%%%%%%%%%%%%%%%%%%%%%
\subsection{MGF of a Bernoulli Distribution}
\label{subsec:mgf-bernoulli}
%%%%%%%%%%%%%%%%%%%%%%%%%%%%%%%%%%%%%%%%%%%%%%%%%%%%%%%%%%%%

If

\[
X\sim\operatorname{Bernoulli}(p),
\]

then

\[
M_X(t)
=
(1-p)e^0+pe^t.
\]

Therefore,

\[
\boxed{
M_X(t)
=
1-p+pe^t.
}
\]

%%%%%%%%%%%%%%%%%%%%%%%%%%%%%%%%%%%%%%%%%%%%%%%%%%%%%%%%%%%%
\subsection{MGF of a Binomial Distribution}
\label{subsec:mgf-binomial}
%%%%%%%%%%%%%%%%%%%%%%%%%%%%%%%%%%%%%%%%%%%%%%%%%%%%%%%%%%%%

If

\[
X\sim\operatorname{Bin}(n,p),
\]

then

\[
\boxed{
M_X(t)
=
(1-p+pe^t)^n.
}
\]

This can be derived either directly or by representing \(X\) as a sum
of independent Bernoulli variables.

%%%%%%%%%%%%%%%%%%%%%%%%%%%%%%%%%%%%%%%%%%%%%%%%%%%%%%%%%%%%
\subsection{MGF of a Poisson Distribution}
\label{subsec:mgf-poisson}
%%%%%%%%%%%%%%%%%%%%%%%%%%%%%%%%%%%%%%%%%%%%%%%%%%%%%%%%%%%%

For

\[
X\sim\operatorname{Poisson}(\lambda),
\]

\[
M_X(t)
=
\sum_{k=0}^{\infty}
e^{tk}
e^{-\lambda}
\frac{\lambda^k}{k!}.
\]

Thus,

\[
M_X(t)
=
e^{-\lambda}
\sum_{k=0}^{\infty}
\frac{
(\lambda e^t)^k
}{
k!
}.
\]

Therefore,

\[
\boxed{
M_X(t)
=
\exp
\left(
\lambda(e^t-1)
\right).
}
\]

%%%%%%%%%%%%%%%%%%%%%%%%%%%%%%%%%%%%%%%%%%%%%%%%%%%%%%%%%%%%
\subsection{MGF of a Normal Distribution}
\label{subsec:mgf-normal}
%%%%%%%%%%%%%%%%%%%%%%%%%%%%%%%%%%%%%%%%%%%%%%%%%%%%%%%%%%%%

If

\[
X\sim N(\mu,\sigma^2),
\]

then

\[
\boxed{
M_X(t)
=
\exp
\left(
\mu t
+
\frac{\sigma^2t^2}{2}
\right).
}
\]

For the standard normal,

\[
\boxed{
M_Z(t)
=
e^{t^2/2}.
}
\]

%%%%%%%%%%%%%%%%%%%%%%%%%%%%%%%%%%%%%%%%%%%%%%%%%%%%%%%%%%%%
\subsection{Worked Example: Normal Moments from the MGF}
\label{subsec:worked-normal-mgf-moments}
%%%%%%%%%%%%%%%%%%%%%%%%%%%%%%%%%%%%%%%%%%%%%%%%%%%%%%%%%%%%

\begin{workedexamplebox}{Recovering Mean and Variance}

Suppose

\[
X\sim N(\mu,\sigma^2).
\]

Its MGF is

\[
M_X(t)
=
\exp
\left(
\mu t+\frac{\sigma^2t^2}{2}
\right).
\]

Differentiate:

\[
M_X'(t)
=
(\mu+\sigma^2t)
M_X(t).
\]

Therefore,

\[
M_X'(0)=\mu.
\]

A second differentiation gives

\[
M_X''(0)
=
\mu^2+\sigma^2.
\]

Hence,

\[
\operatorname{Var}(X)
=
M_X''(0)
-
[M_X'(0)]^2.
\]

\begin{answerbox}
\[
\boxed{
\mathbb{E}[X]=\mu,
\qquad
\operatorname{Var}(X)=\sigma^2.
}
\]
\end{answerbox}

\end{workedexamplebox}

%%%%%%%%%%%%%%%%%%%%%%%%%%%%%%%%%%%%%%%%%%%%%%%%%%%%%%%%%%%%
\subsection{MGF of an Exponential Distribution}
\label{subsec:mgf-exponential}
%%%%%%%%%%%%%%%%%%%%%%%%%%%%%%%%%%%%%%%%%%%%%%%%%%%%%%%%%%%%

If

\[
X\sim\operatorname{Exp}(\lambda),
\]

then

\[
M_X(t)
=
\int_0^\infty
e^{tx}
\lambda e^{-\lambda x}
\,dx.
\]

Thus,

\[
=
\lambda
\int_0^\infty
e^{-(\lambda-t)x}
\,dx.
\]

This converges when

\[
t<\lambda.
\]

Therefore,

\[
\boxed{
M_X(t)
=
\frac{\lambda}{\lambda-t},
\qquad
t<\lambda.
}
\]

%%%%%%%%%%%%%%%%%%%%%%%%%%%%%%%%%%%%%%%%%%%%%%%%%%%%%%%%%%%%
\subsection{MGF of a Gamma Distribution}
\label{subsec:mgf-gamma}
%%%%%%%%%%%%%%%%%%%%%%%%%%%%%%%%%%%%%%%%%%%%%%%%%%%%%%%%%%%%

For

\[
X\sim\operatorname{Gamma}(\alpha,\lambda)
\]

using the shape-rate parameterization,

\[
\boxed{
M_X(t)
=
\left(
\frac{\lambda}{\lambda-t}
\right)^\alpha,
\qquad
t<\lambda.
}
\]

Setting

\[
\alpha=1
\]

recovers the exponential MGF.

%%%%%%%%%%%%%%%%%%%%%%%%%%%%%%%%%%%%%%%%%%%%%%%%%%%%%%%%%%%%
\subsection{MGF of a Chi-Square Distribution}
\label{subsec:mgf-chi-square}
%%%%%%%%%%%%%%%%%%%%%%%%%%%%%%%%%%%%%%%%%%%%%%%%%%%%%%%%%%%%

Since

\[
\chi^2_\nu
=
\operatorname{Gamma}
\left(
\frac{\nu}{2},
\frac12
\right),
\]

its MGF is

\[
\boxed{
M_X(t)
=
(1-2t)^{-\nu/2},
\qquad
t<\frac12.
}
\]

%%%%%%%%%%%%%%%%%%%%%%%%%%%%%%%%%%%%%%%%%%%%%%%%%%%%%%%%%%%%
\subsection{MGFs of Independent Sums}
\label{subsec:mgf-independent-sums}
%%%%%%%%%%%%%%%%%%%%%%%%%%%%%%%%%%%%%%%%%%%%%%%%%%%%%%%%%%%%

If \(X\) and \(Y\) are independent and their MGFs exist, then

\[
M_{X+Y}(t)
=
\mathbb{E}[e^{t(X+Y)}].
\]

Since

\[
e^{t(X+Y)}
=
e^{tX}e^{tY},
\]

independence gives

\[
\boxed{
M_{X+Y}(t)
=
M_X(t)M_Y(t).
}
\]

%%%%%%%%%%%%%%%%%%%%%%%%%%%%%%%%%%%%%%%%%%%%%%%%%%%%%%%%%%%%
\subsection{Worked Example: Sum of Independent Normals}
\label{subsec:worked-mgf-normal-sum}
%%%%%%%%%%%%%%%%%%%%%%%%%%%%%%%%%%%%%%%%%%%%%%%%%%%%%%%%%%%%

\begin{workedexamplebox}{Normal Closure by MGFs}

Suppose

\[
X\sim N(\mu_1,\sigma_1^2)
\]

and

\[
Y\sim N(\mu_2,\sigma_2^2)
\]

independently.

Then

\[
M_{X+Y}(t)
=
M_X(t)M_Y(t).
\]

Therefore,

\[
=
\exp
\left(
\mu_1t+\frac{\sigma_1^2t^2}{2}
\right)
\exp
\left(
\mu_2t+\frac{\sigma_2^2t^2}{2}
\right).
\]

Combining exponents,

\[
=
\exp
\left[
(\mu_1+\mu_2)t
+
\frac{
(\sigma_1^2+\sigma_2^2)t^2
}{2}
\right].
\]

\begin{answerbox}
\[
\boxed{
X+Y
\sim
N
\left(
\mu_1+\mu_2,
\sigma_1^2+\sigma_2^2
\right).
}
\]
\end{answerbox}

\end{workedexamplebox}

%%%%%%%%%%%%%%%%%%%%%%%%%%%%%%%%%%%%%%%%%%%%%%%%%%%%%%%%%%%%
\subsection{MGF Uniqueness}
\label{subsec:mgf-uniqueness}
%%%%%%%%%%%%%%%%%%%%%%%%%%%%%%%%%%%%%%%%%%%%%%%%%%%%%%%%%%%%

Under suitable conditions, if two random variables have the same MGF in
an open interval containing \(0\), then they have the same distribution.

Thus,

\[
\boxed{
M_X(t)=M_Y(t)
\text{ near }0
\Longrightarrow
X\overset{d}{=}Y.
}
\]

This allows transforms to identify distributions.

%%%%%%%%%%%%%%%%%%%%%%%%%%%%%%%%%%%%%%%%%%%%%%%%%%%%%%%%%%%%
\subsection{MGFs Do Not Always Exist}
\label{subsec:mgf-not-always-exist}
%%%%%%%%%%%%%%%%%%%%%%%%%%%%%%%%%%%%%%%%%%%%%%%%%%%%%%%%%%%%

An MGF may fail to exist for nonzero values of \(t\).

Heavy-tailed distributions provide important examples.

For instance, the standard Cauchy distribution does not have a finite
MGF in a neighborhood of zero.

Therefore,

\[
\boxed{
\text{MGFs are powerful but not universal}.
}
\]

This motivates characteristic functions.

%%%%%%%%%%%%%%%%%%%%%%%%%%%%%%%%%%%%%%%%%%%%%%%%%%%%%%%%%%%%
\subsection{Characteristic Functions}
\label{subsec:characteristic-functions}
%%%%%%%%%%%%%%%%%%%%%%%%%%%%%%%%%%%%%%%%%%%%%%%%%%%%%%%%%%%%

\begin{definitionbox}{Characteristic Function}
The characteristic function of a random variable \(X\) is

\[
\boxed{
\varphi_X(t)
=
\mathbb{E}[e^{itX}],
\qquad
t\in\mathbb{R},
}
\]

where

\[
i^2=-1.
\]
\end{definitionbox}

The Greek letter

\[
\varphi
\]

is called ``varphi'' or ``phi.''

%%%%%%%%%%%%%%%%%%%%%%%%%%%%%%%%%%%%%%%%%%%%%%%%%%%%%%%%%%%%
\subsection{Euler's Formula and Characteristic Functions}
\label{subsec:euler-characteristic}
%%%%%%%%%%%%%%%%%%%%%%%%%%%%%%%%%%%%%%%%%%%%%%%%%%%%%%%%%%%%

Euler's formula gives

\[
e^{itX}
=
\cos(tX)
+
i\sin(tX).
\]

Therefore,

\[
\boxed{
\varphi_X(t)
=
\mathbb{E}[\cos(tX)]
+
i\mathbb{E}[\sin(tX)].
}
\]

Thus a characteristic function combines cosine and sine transforms of
the distribution.

%%%%%%%%%%%%%%%%%%%%%%%%%%%%%%%%%%%%%%%%%%%%%%%%%%%%%%%%%%%%
\subsection{Characteristic Functions Always Exist}
\label{subsec:characteristic-functions-exist}
%%%%%%%%%%%%%%%%%%%%%%%%%%%%%%%%%%%%%%%%%%%%%%%%%%%%%%%%%%%%

Since

\[
|e^{itX}|=1,
\]

we have

\[
\left|
\mathbb{E}[e^{itX}]
\right|
\leq1.
\]

Therefore,

\[
\boxed{
|\varphi_X(t)|\leq1.
}
\]

Unlike MGFs, characteristic functions exist for every real-valued random
variable.

%%%%%%%%%%%%%%%%%%%%%%%%%%%%%%%%%%%%%%%%%%%%%%%%%%%%%%%%%%%%
\subsection{Basic Characteristic Function Properties}
\label{subsec:characteristic-basic-properties}
%%%%%%%%%%%%%%%%%%%%%%%%%%%%%%%%%%%%%%%%%%%%%%%%%%%%%%%%%%%%

Every characteristic function satisfies

\[
\boxed{
\varphi_X(0)=1.
}
\]

Also,

\[
\boxed{
|\varphi_X(t)|\leq1.
}
\]

Furthermore,

\[
\boxed{
\varphi_X(-t)
=
\overline{\varphi_X(t)},
}
\]

where the overline denotes complex conjugation.

%%%%%%%%%%%%%%%%%%%%%%%%%%%%%%%%%%%%%%%%%%%%%%%%%%%%%%%%%%%%
\subsection{Characteristic Function of a Bernoulli Variable}
\label{subsec:cf-bernoulli}
%%%%%%%%%%%%%%%%%%%%%%%%%%%%%%%%%%%%%%%%%%%%%%%%%%%%%%%%%%%%

If

\[
X\sim\operatorname{Bernoulli}(p),
\]

then

\[
\varphi_X(t)
=
(1-p)+pe^{it}.
\]

Therefore,

\[
\boxed{
\varphi_X(t)
=
1-p+pe^{it}.
}
\]

%%%%%%%%%%%%%%%%%%%%%%%%%%%%%%%%%%%%%%%%%%%%%%%%%%%%%%%%%%%%
\subsection{Characteristic Function of a Binomial Variable}
\label{subsec:cf-binomial}
%%%%%%%%%%%%%%%%%%%%%%%%%%%%%%%%%%%%%%%%%%%%%%%%%%%%%%%%%%%%

If

\[
X\sim\operatorname{Bin}(n,p),
\]

then

\[
\boxed{
\varphi_X(t)
=
(1-p+pe^{it})^n.
}
\]

This is obtained from the MGF by formally replacing \(t\) with \(it\).

%%%%%%%%%%%%%%%%%%%%%%%%%%%%%%%%%%%%%%%%%%%%%%%%%%%%%%%%%%%%
\subsection{Characteristic Function of a Poisson Variable}
\label{subsec:cf-poisson}
%%%%%%%%%%%%%%%%%%%%%%%%%%%%%%%%%%%%%%%%%%%%%%%%%%%%%%%%%%%%

For

\[
X\sim\operatorname{Poisson}(\lambda),
\]

\[
\boxed{
\varphi_X(t)
=
\exp
\left(
\lambda(e^{it}-1)
\right).
}
\]

%%%%%%%%%%%%%%%%%%%%%%%%%%%%%%%%%%%%%%%%%%%%%%%%%%%%%%%%%%%%
\subsection{Characteristic Function of a Normal Variable}
\label{subsec:cf-normal}
%%%%%%%%%%%%%%%%%%%%%%%%%%%%%%%%%%%%%%%%%%%%%%%%%%%%%%%%%%%%

If

\[
X\sim N(\mu,\sigma^2),
\]

then

\[
\boxed{
\varphi_X(t)
=
\exp
\left(
i\mu t
-
\frac{\sigma^2t^2}{2}
\right).
}
\]

For the standard normal,

\[
\boxed{
\varphi_Z(t)
=
e^{-t^2/2}.
}
\]

%%%%%%%%%%%%%%%%%%%%%%%%%%%%%%%%%%%%%%%%%%%%%%%%%%%%%%%%%%%%
\subsection{Characteristic Function of a Cauchy Variable}
\label{subsec:cf-cauchy}
%%%%%%%%%%%%%%%%%%%%%%%%%%%%%%%%%%%%%%%%%%%%%%%%%%%%%%%%%%%%

For a standard Cauchy random variable,

\[
\boxed{
\varphi_X(t)
=
e^{-|t|}.
}
\]

This characteristic function exists even though the Cauchy MGF does
not.

This illustrates one major advantage of characteristic functions.

%%%%%%%%%%%%%%%%%%%%%%%%%%%%%%%%%%%%%%%%%%%%%%%%%%%%%%%%%%%%
\subsection{Characteristic Functions of Independent Sums}
\label{subsec:cf-independent-sums}
%%%%%%%%%%%%%%%%%%%%%%%%%%%%%%%%%%%%%%%%%%%%%%%%%%%%%%%%%%%%

If \(X\) and \(Y\) are independent, then

\[
\varphi_{X+Y}(t)
=
\mathbb{E}
[
e^{it(X+Y)}
].
\]

Since

\[
e^{it(X+Y)}
=
e^{itX}e^{itY},
\]

independence gives

\[
\boxed{
\varphi_{X+Y}(t)
=
\varphi_X(t)\varphi_Y(t).
}
\]

Thus characteristic functions convert convolution into multiplication.

%%%%%%%%%%%%%%%%%%%%%%%%%%%%%%%%%%%%%%%%%%%%%%%%%%%%%%%%%%%%
\subsection{Characteristic Function Uniqueness}
\label{subsec:cf-uniqueness}
%%%%%%%%%%%%%%%%%%%%%%%%%%%%%%%%%%%%%%%%%%%%%%%%%%%%%%%%%%%%

Characteristic functions uniquely determine probability distributions.

Thus,

\[
\boxed{
\varphi_X(t)
=
\varphi_Y(t)
\text{ for all }t
\Longrightarrow
X\overset{d}{=}Y.
}
\]

Unlike MGF uniqueness, this result does not require existence in a
neighborhood of zero because characteristic functions always exist.

%%%%%%%%%%%%%%%%%%%%%%%%%%%%%%%%%%%%%%%%%%%%%%%%%%%%%%%%%%%%
\subsection{Characteristic Functions and Moments}
\label{subsec:cf-moments}
%%%%%%%%%%%%%%%%%%%%%%%%%%%%%%%%%%%%%%%%%%%%%%%%%%%%%%%%%%%%

If sufficiently many moments exist, derivatives of the characteristic
function recover them.

For example,

\[
\varphi_X'(0)
=
i\mathbb{E}[X].
\]

More generally,

\[
\boxed{
\varphi_X^{(k)}(0)
=
i^k
\mathbb{E}[X^k].
}
\]

Hence,

\[
\boxed{
\mathbb{E}[X^k]
=
\frac{
\varphi_X^{(k)}(0)
}{
i^k
}.
}
\]

%%%%%%%%%%%%%%%%%%%%%%%%%%%%%%%%%%%%%%%%%%%%%%%%%%%%%%%%%%%%
\subsection{Characteristic Functions and Fourier Analysis}
\label{subsec:cf-fourier-analysis}
%%%%%%%%%%%%%%%%%%%%%%%%%%%%%%%%%%%%%%%%%%%%%%%%%%%%%%%%%%%%

For a continuous random variable with density \(f_X\),

\[
\varphi_X(t)
=
\int_{-\infty}^{\infty}
e^{itx}
f_X(x)\,dx.
\]

Thus the characteristic function is essentially a Fourier transform of
the probability density.

Therefore,

\[
\boxed{
\text{characteristic functions}
\longleftrightarrow
\text{Fourier analysis}.
}
\]

%%%%%%%%%%%%%%%%%%%%%%%%%%%%%%%%%%%%%%%%%%%%%%%%%%%%%%%%%%%%
\subsection{Fourier Inversion Preview}
\label{subsec:fourier-inversion-probability}
%%%%%%%%%%%%%%%%%%%%%%%%%%%%%%%%%%%%%%%%%%%%%%%%%%%%%%%%%%%%

Under appropriate regularity conditions, the density can be recovered
from the characteristic function through Fourier inversion:

\[
\boxed{
f_X(x)
=
\frac1{2\pi}
\int_{-\infty}^{\infty}
e^{-itx}
\varphi_X(t)\,dt.
}
\]

Thus transform and distribution can be viewed as two representations of
the same probability law.

%%%%%%%%%%%%%%%%%%%%%%%%%%%%%%%%%%%%%%%%%%%%%%%%%%%%%%%%%%%%
\subsection{Characteristic Functions and Convergence}
\label{subsec:cf-convergence-preview}
%%%%%%%%%%%%%%%%%%%%%%%%%%%%%%%%%%%%%%%%%%%%%%%%%%%%%%%%%%%%

Characteristic functions are especially important in limit theorems.

A major result states, roughly, that convergence of characteristic
functions to a valid limiting characteristic function corresponds to
convergence in distribution.

This principle is formalized by Lévy's continuity theorem.

It will be used conceptually in the study of central limit theorems.

%%%%%%%%%%%%%%%%%%%%%%%%%%%%%%%%%%%%%%%%%%%%%%%%%%%%%%%%%%%%
\subsection{Lévy Continuity Theorem Preview}
\label{subsec:levy-continuity-preview}
%%%%%%%%%%%%%%%%%%%%%%%%%%%%%%%%%%%%%%%%%%%%%%%%%%%%%%%%%%%%

Suppose

\[
\varphi_{X_n}(t)
\to
\varphi(t)
\]

for every real \(t\), and suppose \(\varphi\) is continuous at \(0\) and
is the characteristic function of some random variable \(X\).

Then, under the standard theorem,

\[
\boxed{
X_n
\xrightarrow{d}
X.
}
\]

The symbol

\[
\xrightarrow{d}
\]

means ``converges in distribution.''

%%%%%%%%%%%%%%%%%%%%%%%%%%%%%%%%%%%%%%%%%%%%%%%%%%%%%%%%%%%%
\subsection{Laplace Transforms}
\label{subsec:laplace-transforms-probability}
%%%%%%%%%%%%%%%%%%%%%%%%%%%%%%%%%%%%%%%%%%%%%%%%%%%%%%%%%%%%

For nonnegative random variables, Laplace transforms are often more
natural than MGFs.

\begin{definitionbox}{Laplace Transform of a Random Variable}
For a nonnegative random variable \(X\), define

\[
\boxed{
L_X(s)
=
\mathbb{E}[e^{-sX}],
\qquad
s\geq0.
}
\]
\end{definitionbox}

For a continuous variable,

\[
\boxed{
L_X(s)
=
\int_0^\infty
e^{-sx}f_X(x)\,dx.
}
\]

%%%%%%%%%%%%%%%%%%%%%%%%%%%%%%%%%%%%%%%%%%%%%%%%%%%%%%%%%%%%
\subsection{Relation Between Laplace Transform and MGF}
\label{subsec:laplace-mgf-relation}
%%%%%%%%%%%%%%%%%%%%%%%%%%%%%%%%%%%%%%%%%%%%%%%%%%%%%%%%%%%%

Whenever both expressions are defined,

\[
M_X(t)
=
\mathbb{E}[e^{tX}]
\]

and

\[
L_X(s)
=
\mathbb{E}[e^{-sX}].
\]

Therefore,

\[
\boxed{
L_X(s)
=
M_X(-s).
}
\]

Thus the Laplace transform is the MGF evaluated at a negative argument.

%%%%%%%%%%%%%%%%%%%%%%%%%%%%%%%%%%%%%%%%%%%%%%%%%%%%%%%%%%%%
\subsection{Laplace Transform of an Exponential Variable}
\label{subsec:laplace-exponential}
%%%%%%%%%%%%%%%%%%%%%%%%%%%%%%%%%%%%%%%%%%%%%%%%%%%%%%%%%%%%

If

\[
X\sim\operatorname{Exp}(\lambda),
\]

then

\[
L_X(s)
=
\int_0^\infty
e^{-sx}
\lambda e^{-\lambda x}
\,dx.
\]

Thus,

\[
=
\lambda
\int_0^\infty
e^{-(\lambda+s)x}
\,dx.
\]

Therefore,

\[
\boxed{
L_X(s)
=
\frac{\lambda}{\lambda+s}.
}
\]

%%%%%%%%%%%%%%%%%%%%%%%%%%%%%%%%%%%%%%%%%%%%%%%%%%%%%%%%%%%%
\subsection{Laplace Transform of a Gamma Variable}
\label{subsec:laplace-gamma}
%%%%%%%%%%%%%%%%%%%%%%%%%%%%%%%%%%%%%%%%%%%%%%%%%%%%%%%%%%%%

If

\[
X\sim\operatorname{Gamma}(\alpha,\lambda)
\]

under the shape-rate convention, then

\[
\boxed{
L_X(s)
=
\left(
\frac{\lambda}{\lambda+s}
\right)^\alpha.
}
\]

%%%%%%%%%%%%%%%%%%%%%%%%%%%%%%%%%%%%%%%%%%%%%%%%%%%%%%%%%%%%
\subsection{Laplace Transforms of Independent Sums}
\label{subsec:laplace-independent-sums}
%%%%%%%%%%%%%%%%%%%%%%%%%%%%%%%%%%%%%%%%%%%%%%%%%%%%%%%%%%%%

If \(X\) and \(Y\) are independent nonnegative random variables, then

\[
L_{X+Y}(s)
=
\mathbb{E}[e^{-s(X+Y)}].
\]

Independence gives

\[
\boxed{
L_{X+Y}(s)
=
L_X(s)L_Y(s).
}
\]

Again, convolution becomes multiplication.

%%%%%%%%%%%%%%%%%%%%%%%%%%%%%%%%%%%%%%%%%%%%%%%%%%%%%%%%%%%%
\subsection{Worked Example: Gamma Closure}
\label{subsec:worked-laplace-gamma-sum}
%%%%%%%%%%%%%%%%%%%%%%%%%%%%%%%%%%%%%%%%%%%%%%%%%%%%%%%%%%%%

\begin{workedexamplebox}{Sum of Independent Gamma Variables}

Suppose

\[
X\sim\operatorname{Gamma}(\alpha_1,\lambda)
\]

and

\[
Y\sim\operatorname{Gamma}(\alpha_2,\lambda)
\]

independently.

Then

\[
L_X(s)
=
\left(
\frac{\lambda}{\lambda+s}
\right)^{\alpha_1}
\]

and

\[
L_Y(s)
=
\left(
\frac{\lambda}{\lambda+s}
\right)^{\alpha_2}.
\]

Therefore,

\[
L_{X+Y}(s)
=
\left(
\frac{\lambda}{\lambda+s}
\right)^{\alpha_1+\alpha_2}.
\]

\begin{answerbox}
\[
\boxed{
X+Y
\sim
\operatorname{Gamma}
(\alpha_1+\alpha_2,\lambda).
}
\]
\end{answerbox}

\end{workedexamplebox}

%%%%%%%%%%%%%%%%%%%%%%%%%%%%%%%%%%%%%%%%%%%%%%%%%%%%%%%%%%%%
\subsection{Cumulant Generating Functions}
\label{subsec:cumulant-generating-functions}
%%%%%%%%%%%%%%%%%%%%%%%%%%%%%%%%%%%%%%%%%%%%%%%%%%%%%%%%%%%%

The logarithm of the MGF produces another important transform.

\begin{definitionbox}{Cumulant Generating Function}
If \(M_X(t)\) exists, the cumulant generating function is

\[
\boxed{
K_X(t)
=
\ln M_X(t).
}
\]
\end{definitionbox}

Its derivatives at zero are called cumulants.

%%%%%%%%%%%%%%%%%%%%%%%%%%%%%%%%%%%%%%%%%%%%%%%%%%%%%%%%%%%%
\subsection{Cumulants}
\label{subsec:cumulants}
%%%%%%%%%%%%%%%%%%%%%%%%%%%%%%%%%%%%%%%%%%%%%%%%%%%%%%%%%%%%

The \(k\)-th cumulant is

\[
\boxed{
\kappa_k
=
K_X^{(k)}(0),
}
\]

when the derivative exists.

The first few cumulants have familiar meanings.

%%%%%%%%%%%%%%%%%%%%%%%%%%%%%%%%%%%%%%%%%%%%%%%%%%%%%%%%%%%%
\subsection{First Cumulant}
\label{subsec:first-cumulant}
%%%%%%%%%%%%%%%%%%%%%%%%%%%%%%%%%%%%%%%%%%%%%%%%%%%%%%%%%%%%

Since

\[
K_X(t)
=
\ln M_X(t),
\]

we have

\[
K_X'(t)
=
\frac{
M_X'(t)
}{
M_X(t)
}.
\]

At \(t=0\),

\[
M_X(0)=1,
\]

so

\[
\boxed{
\kappa_1
=
\mathbb{E}[X].
}
\]

Thus the first cumulant is the mean.

%%%%%%%%%%%%%%%%%%%%%%%%%%%%%%%%%%%%%%%%%%%%%%%%%%%%%%%%%%%%
\subsection{Second Cumulant}
\label{subsec:second-cumulant}
%%%%%%%%%%%%%%%%%%%%%%%%%%%%%%%%%%%%%%%%%%%%%%%%%%%%%%%%%%%%

Differentiating again yields

\[
\boxed{
\kappa_2
=
\operatorname{Var}(X).
}
\]

Thus the second cumulant is the variance.

%%%%%%%%%%%%%%%%%%%%%%%%%%%%%%%%%%%%%%%%%%%%%%%%%%%%%%%%%%%%
\subsection{Third and Fourth Cumulants}
\label{subsec:higher-cumulants}
%%%%%%%%%%%%%%%%%%%%%%%%%%%%%%%%%%%%%%%%%%%%%%%%%%%%%%%%%%%%

The third cumulant is the third central moment:

\[
\boxed{
\kappa_3
=
\mathbb{E}
[
(X-\mu)^3
].
}
\]

The fourth cumulant is

\[
\boxed{
\kappa_4
=
\mathbb{E}
[
(X-\mu)^4
]
-
3\sigma^4.
}
\]

Thus the fourth cumulant measures excess fourth-order behavior relative
to the normal distribution.

%%%%%%%%%%%%%%%%%%%%%%%%%%%%%%%%%%%%%%%%%%%%%%%%%%%%%%%%%%%%
\subsection{Cumulants of Independent Sums}
\label{subsec:cumulants-independent-sums}
%%%%%%%%%%%%%%%%%%%%%%%%%%%%%%%%%%%%%%%%%%%%%%%%%%%%%%%%%%%%

If \(X\) and \(Y\) are independent, then

\[
M_{X+Y}(t)
=
M_X(t)M_Y(t).
\]

Taking logarithms,

\[
K_{X+Y}(t)
=
K_X(t)+K_Y(t).
\]

Therefore,

\[
\boxed{
\kappa_k(X+Y)
=
\kappa_k(X)
+
\kappa_k(Y).
}
\]

Cumulants add under independent sums.

%%%%%%%%%%%%%%%%%%%%%%%%%%%%%%%%%%%%%%%%%%%%%%%%%%%%%%%%%%%%
\subsection{Worked Example: Poisson Cumulants}
\label{subsec:worked-poisson-cumulants}
%%%%%%%%%%%%%%%%%%%%%%%%%%%%%%%%%%%%%%%%%%%%%%%%%%%%%%%%%%%%

\begin{workedexamplebox}{All Poisson Cumulants}

For

\[
X\sim\operatorname{Poisson}(\lambda),
\]

the MGF is

\[
M_X(t)
=
\exp
\left(
\lambda(e^t-1)
\right).
\]

Therefore,

\[
K_X(t)
=
\lambda(e^t-1).
\]

Every derivative of \(K_X\) at \(0\) equals \(\lambda\).

Thus,

\begin{answerbox}
\[
\boxed{
\kappa_k
=
\lambda
}
\]

for every positive integer \(k\).
\end{answerbox}

\end{workedexamplebox}

%%%%%%%%%%%%%%%%%%%%%%%%%%%%%%%%%%%%%%%%%%%%%%%%%%%%%%%%%%%%
\subsection{Cumulants of a Normal Distribution}
\label{subsec:normal-cumulants}
%%%%%%%%%%%%%%%%%%%%%%%%%%%%%%%%%%%%%%%%%%%%%%%%%%%%%%%%%%%%

For

\[
X\sim N(\mu,\sigma^2),
\]

\[
M_X(t)
=
\exp
\left(
\mu t+\frac{\sigma^2t^2}{2}
\right).
\]

Therefore,

\[
K_X(t)
=
\mu t
+
\frac{\sigma^2t^2}{2}.
\]

Hence,

\[
\boxed{
\kappa_1=\mu,
}
\]

\[
\boxed{
\kappa_2=\sigma^2,
}
\]

and

\[
\boxed{
\kappa_k=0
\quad
\text{for }k\geq3.
}
\]

The vanishing of all higher cumulants is a defining structural feature
of the normal distribution.

%%%%%%%%%%%%%%%%%%%%%%%%%%%%%%%%%%%%%%%%%%%%%%%%%%%%%%%%%%%%
\subsection{Scaling Cumulants}
\label{subsec:scaling-cumulants}
%%%%%%%%%%%%%%%%%%%%%%%%%%%%%%%%%%%%%%%%%%%%%%%%%%%%%%%%%%%%

For

\[
Y=aX+b,
\]

the cumulants transform as

\[
\boxed{
\kappa_1(Y)
=
a\kappa_1(X)+b,
}
\]

while for

\[
k\geq2,
\]

\[
\boxed{
\kappa_k(Y)
=
a^k\kappa_k(X).
}
\]

Thus translation affects only the first cumulant.

%%%%%%%%%%%%%%%%%%%%%%%%%%%%%%%%%%%%%%%%%%%%%%%%%%%%%%%%%%%%
\subsection{Transforms and Convolution}
\label{subsec:transforms-and-convolution}
%%%%%%%%%%%%%%%%%%%%%%%%%%%%%%%%%%%%%%%%%%%%%%%%%%%%%%%%%%%%

Convolution appears naturally when independent variables are added.

In the original distribution domain:

\[
\boxed{
f_{X+Y}
=
f_X*f_Y.
}
\]

In the transform domain:

\[
\boxed{
T_{X+Y}
=
T_XT_Y.
}
\]

Therefore transforms convert a convolution problem into multiplication.

This is one of their most important advantages.

%%%%%%%%%%%%%%%%%%%%%%%%%%%%%%%%%%%%%%%%%%%%%%%%%%%%%%%%%%%%
\subsection{Transform Table}
\label{subsec:probability-transform-table}
%%%%%%%%%%%%%%%%%%%%%%%%%%%%%%%%%%%%%%%%%%%%%%%%%%%%%%%%%%%%

\[
\begin{array}{c|c|c}
\text{Transform}
&
\text{Definition}
&
\text{Typical Setting}
\\
\hline
G_X(s)
&
\mathbb{E}[s^X]
&
\text{nonnegative integer-valued }X
\\[1mm]
M_X(t)
&
\mathbb{E}[e^{tX}]
&
\text{moments, sums}
\\[1mm]
\varphi_X(t)
&
\mathbb{E}[e^{itX}]
&
\text{general distributions}
\\[1mm]
L_X(s)
&
\mathbb{E}[e^{-sX}]
&
\text{nonnegative }X
\\[1mm]
K_X(t)
&
\ln M_X(t)
&
\text{cumulants}
\end{array}
\]

%%%%%%%%%%%%%%%%%%%%%%%%%%%%%%%%%%%%%%%%%%%%%%%%%%%%%%%%%%%%
\subsection{Relationship Between PGF and MGF}
\label{subsec:pgf-mgf-relation}
%%%%%%%%%%%%%%%%%%%%%%%%%%%%%%%%%%%%%%%%%%%%%%%%%%%%%%%%%%%%

For a nonnegative integer-valued \(X\),

\[
M_X(t)
=
\mathbb{E}[e^{tX}].
\]

Since

\[
e^{tX}
=
(e^t)^X,
\]

we obtain

\[
\boxed{
M_X(t)
=
G_X(e^t).
}
\]

Thus the MGF can be obtained from the PGF by substituting

\[
s=e^t.
\]

%%%%%%%%%%%%%%%%%%%%%%%%%%%%%%%%%%%%%%%%%%%%%%%%%%%%%%%%%%%%
\subsection{Relationship Between MGF and Characteristic Function}
\label{subsec:mgf-cf-relation}
%%%%%%%%%%%%%%%%%%%%%%%%%%%%%%%%%%%%%%%%%%%%%%%%%%%%%%%%%%%%

Whenever the MGF has an appropriate complex extension,

\[
\boxed{
\varphi_X(t)
=
M_X(it).
}
\]

The characteristic function replaces the real exponential factor with
a complex oscillatory factor.

%%%%%%%%%%%%%%%%%%%%%%%%%%%%%%%%%%%%%%%%%%%%%%%%%%%%%%%%%%%%
\subsection{Transforms of Linear Transformations}
\label{subsec:transforms-linear-transformations}
%%%%%%%%%%%%%%%%%%%%%%%%%%%%%%%%%%%%%%%%%%%%%%%%%%%%%%%%%%%%

If

\[
Y=aX+b,
\]

then

\[
M_Y(t)
=
\mathbb{E}
[
e^{t(aX+b)}
].
\]

Therefore,

\[
\boxed{
M_Y(t)
=
e^{bt}
M_X(at).
}
\]

Similarly,

\[
\boxed{
\varphi_Y(t)
=
e^{ibt}
\varphi_X(at).
}
\]

%%%%%%%%%%%%%%%%%%%%%%%%%%%%%%%%%%%%%%%%%%%%%%%%%%%%%%%%%%%%
\subsection{Worked Example: Standardization via Characteristic Functions}
\label{subsec:worked-cf-standardization}
%%%%%%%%%%%%%%%%%%%%%%%%%%%%%%%%%%%%%%%%%%%%%%%%%%%%%%%%%%%%

\begin{workedexamplebox}{Standardizing a Normal Variable}

Suppose

\[
X\sim N(\mu,\sigma^2)
\]

and define

\[
Z
=
\frac{X-\mu}{\sigma}.
\]

Using the linear-transform rule,

\[
\varphi_Z(t)
=
e^{-it\mu/\sigma}
\varphi_X(t/\sigma).
\]

Since

\[
\varphi_X(u)
=
\exp
\left(
i\mu u
-
\frac{\sigma^2u^2}{2}
\right),
\]

substituting

\[
u=\frac{t}{\sigma}
\]

gives

\[
\varphi_Z(t)
=
e^{-t^2/2}.
\]

\begin{answerbox}
\[
\boxed{
Z\sim N(0,1).
}
\]
\end{answerbox}

\end{workedexamplebox}

%%%%%%%%%%%%%%%%%%%%%%%%%%%%%%%%%%%%%%%%%%%%%%%%%%%%%%%%%%%%
\subsection{Transforms and Distribution Identification}
\label{subsec:transform-distribution-identification}
%%%%%%%%%%%%%%%%%%%%%%%%%%%%%%%%%%%%%%%%%%%%%%%%%%%%%%%%%%%%

A common strategy is:

\begin{enumerate}
    \item compute the transform of a random variable;
    \item simplify the result;
    \item recognize it as the transform of a known distribution;
    \item use uniqueness to identify the distribution.
\end{enumerate}

This approach is especially useful for sums of independent variables.

%%%%%%%%%%%%%%%%%%%%%%%%%%%%%%%%%%%%%%%%%%%%%%%%%%%%%%%%%%%%
\subsection{Worked Example: Chi-Square Additivity}
\label{subsec:worked-chi-square-transform}
%%%%%%%%%%%%%%%%%%%%%%%%%%%%%%%%%%%%%%%%%%%%%%%%%%%%%%%%%%%%

\begin{workedexamplebox}{Adding Independent Chi-Square Variables}

Suppose

\[
X\sim\chi^2_{\nu_1}
\]

and

\[
Y\sim\chi^2_{\nu_2}
\]

independently.

Their MGFs are

\[
M_X(t)
=
(1-2t)^{-\nu_1/2}
\]

and

\[
M_Y(t)
=
(1-2t)^{-\nu_2/2}.
\]

Therefore,

\[
M_{X+Y}(t)
=
(1-2t)^{-(\nu_1+\nu_2)/2}.
\]

\begin{answerbox}
\[
\boxed{
X+Y
\sim
\chi^2_{\nu_1+\nu_2}.
}
\]
\end{answerbox}

\end{workedexamplebox}

%%%%%%%%%%%%%%%%%%%%%%%%%%%%%%%%%%%%%%%%%%%%%%%%%%%%%%%%%%%%
\subsection{Transform Inversion}
\label{subsec:transform-inversion}
%%%%%%%%%%%%%%%%%%%%%%%%%%%%%%%%%%%%%%%%%%%%%%%%%%%%%%%%%%%%

A transform compresses a probability distribution into another
functional representation.

To recover the original distribution, an inversion procedure may be
used.

Examples include:

\[
\boxed{
\text{coefficient extraction for PGFs},
}
\]

\[
\boxed{
\text{Fourier inversion for characteristic functions},
}
\]

and

\[
\boxed{
\text{Laplace inversion for Laplace transforms}.
}
\]

Thus probability transforms are reversible under appropriate
conditions.

%%%%%%%%%%%%%%%%%%%%%%%%%%%%%%%%%%%%%%%%%%%%%%%%%%%%%%%%%%%%
\subsection{Transforms and Recurrence Relations}
\label{subsec:transforms-recurrence-relations}
%%%%%%%%%%%%%%%%%%%%%%%%%%%%%%%%%%%%%%%%%%%%%%%%%%%%%%%%%%%%

PGFs are also useful when probabilities satisfy recurrence relations.

Suppose a sequence

\[
p_0,p_1,p_2,\ldots
\]

satisfies a recurrence.

Multiplying by powers of \(s\) and summing may convert the recurrence
into an algebraic equation for

\[
G(s).
\]

This is the same principle used for generating functions in
combinatorics.

%%%%%%%%%%%%%%%%%%%%%%%%%%%%%%%%%%%%%%%%%%%%%%%%%%%%%%%%%%%%
\subsection{Transforms and Random Sums}
\label{subsec:transforms-random-sums}
%%%%%%%%%%%%%%%%%%%%%%%%%%%%%%%%%%%%%%%%%%%%%%%%%%%%%%%%%%%%

Suppose

\[
N
\]

is a nonnegative integer-valued random variable and

\[
X_1,X_2,\ldots
\]

are i.i.d. random variables independent of \(N\).

Define the random sum

\[
S_N
=
X_1+\cdots+X_N.
\]

Transforms can encode the distribution of this compound random quantity.

%%%%%%%%%%%%%%%%%%%%%%%%%%%%%%%%%%%%%%%%%%%%%%%%%%%%%%%%%%%%
\subsection{PGF of a Random Sum}
\label{subsec:pgf-random-sum}
%%%%%%%%%%%%%%%%%%%%%%%%%%%%%%%%%%%%%%%%%%%%%%%%%%%%%%%%%%%%

Suppose the \(X_i\) are nonnegative integer-valued with common PGF

\[
G_X(s).
\]

Conditioning on \(N=n\),

\[
G_{S_N}(s\mid N=n)
=
[G_X(s)]^n.
\]

Averaging over \(N\),

\[
\boxed{
G_{S_N}(s)
=
G_N(G_X(s)).
}
\]

Thus random summation becomes composition of generating functions.

%%%%%%%%%%%%%%%%%%%%%%%%%%%%%%%%%%%%%%%%%%%%%%%%%%%%%%%%%%%%
\subsection{Compound Poisson Distribution Preview}
\label{subsec:compound-poisson-transform-preview}
%%%%%%%%%%%%%%%%%%%%%%%%%%%%%%%%%%%%%%%%%%%%%%%%%%%%%%%%%%%%

Suppose

\[
N\sim\operatorname{Poisson}(\lambda)
\]

and

\[
S_N
=
X_1+\cdots+X_N.
\]

Then

\[
G_N(s)
=
e^{\lambda(s-1)}.
\]

Therefore,

\[
\boxed{
G_{S_N}(s)
=
\exp
\left(
\lambda(G_X(s)-1)
\right).
}
\]

This is the PGF of a compound Poisson distribution.

%%%%%%%%%%%%%%%%%%%%%%%%%%%%%%%%%%%%%%%%%%%%%%%%%%%%%%%%%%%%
\subsection{Transforms and Branching Processes Preview}
\label{subsec:transforms-branching-processes}
%%%%%%%%%%%%%%%%%%%%%%%%%%%%%%%%%%%%%%%%%%%%%%%%%%%%%%%%%%%%

PGFs are fundamental in branching processes.

If each individual produces a random number of offspring with PGF

\[
G(s),
\]

then composition of \(G\) describes offspring counts across
generations.

For example, the PGF after two generations often involves

\[
\boxed{
G(G(s)).
}
\]

Thus functional iteration encodes population evolution.

%%%%%%%%%%%%%%%%%%%%%%%%%%%%%%%%%%%%%%%%%%%%%%%%%%%%%%%%%%%%
\subsection{Transforms and Queueing Theory Preview}
\label{subsec:transforms-queueing-preview}
%%%%%%%%%%%%%%%%%%%%%%%%%%%%%%%%%%%%%%%%%%%%%%%%%%%%%%%%%%%%

Laplace transforms and PGFs occur frequently in queueing theory.

They are used to describe:

\[
\text{waiting times},
\]

\[
\text{service times},
\]

\[
\text{queue lengths},
\]

and

\[
\text{arrival counts}.
\]

Transform equations can sometimes replace difficult integral or
recurrence equations.

%%%%%%%%%%%%%%%%%%%%%%%%%%%%%%%%%%%%%%%%%%%%%%%%%%%%%%%%%%%%
\subsection{Transforms and Reliability Preview}
\label{subsec:transforms-reliability-preview}
%%%%%%%%%%%%%%%%%%%%%%%%%%%%%%%%%%%%%%%%%%%%%%%%%%%%%%%%%%%%

Laplace transforms are useful for sums of independent lifetime and
waiting-time variables.

For example, if a system lifetime is a sum of independent stages,

\[
T=T_1+\cdots+T_n,
\]

then

\[
\boxed{
L_T(s)
=
\prod_{i=1}^{n}
L_{T_i}(s).
}
\]

This may be easier to analyze than repeated convolution.

%%%%%%%%%%%%%%%%%%%%%%%%%%%%%%%%%%%%%%%%%%%%%%%%%%%%%%%%%%%%
\subsection{Transforms and Limit Theorems}
\label{subsec:transforms-limit-theorems}
%%%%%%%%%%%%%%%%%%%%%%%%%%%%%%%%%%%%%%%%%%%%%%%%%%%%%%%%%%%%

Probability transforms are especially powerful when analyzing long
sums.

Suppose

\[
S_n=X_1+\cdots+X_n.
\]

For independent variables,

\[
\varphi_{S_n}(t)
=
\prod_{j=1}^{n}
\varphi_{X_j}(t).
\]

Products become easier to analyze after logarithms:

\[
\ln\varphi_{S_n}(t)
=
\sum_{j=1}^{n}
\ln\varphi_{X_j}(t).
\]

This is one reason cumulants and characteristic functions are natural
tools for central limit theory.

%%%%%%%%%%%%%%%%%%%%%%%%%%%%%%%%%%%%%%%%%%%%%%%%%%%%%%%%%%%%
\subsection{Small-\(t\) Expansion of a Characteristic Function}
\label{subsec:small-t-characteristic-expansion}
%%%%%%%%%%%%%%%%%%%%%%%%%%%%%%%%%%%%%%%%%%%%%%%%%%%%%%%%%%%%

If \(X\) has finite second moment, then near \(t=0\),

\[
\varphi_X(t)
=
1
+
it\mathbb{E}[X]
-
\frac{t^2}{2}
\mathbb{E}[X^2]
+
o(t^2).
\]

For a centered variable with

\[
\mathbb{E}[X]=0,
\]

this becomes

\[
\boxed{
\varphi_X(t)
=
1
-
\frac{\sigma^2t^2}{2}
+
o(t^2).
}
\]

This expansion underlies transform proofs of the central limit theorem.

%%%%%%%%%%%%%%%%%%%%%%%%%%%%%%%%%%%%%%%%%%%%%%%%%%%%%%%%%%%%
\subsection{Moments Versus Transforms}
\label{subsec:moments-versus-transforms}
%%%%%%%%%%%%%%%%%%%%%%%%%%%%%%%%%%%%%%%%%%%%%%%%%%%%%%%%%%%%

A sequence of moments may describe a distribution, but not every
distribution is uniquely determined by its moments without additional
conditions.

Transforms often encode more complete information.

In particular,

\[
\boxed{
\text{characteristic function}
}
\]

always uniquely determines the probability distribution.

This makes characteristic functions more robust than relying only on
moments.

%%%%%%%%%%%%%%%%%%%%%%%%%%%%%%%%%%%%%%%%%%%%%%%%%%%%%%%%%%%%
\subsection{Moment Determinacy Preview}
\label{subsec:moment-determinacy-preview}
%%%%%%%%%%%%%%%%%%%%%%%%%%%%%%%%%%%%%%%%%%%%%%%%%%%%%%%%%%%%

A distribution is called moment-determinate if its sequence of moments
uniquely determines the distribution.

Some distributions are moment-indeterminate.

Thus,

\[
\boxed{
\text{same moments}
\not\Longrightarrow
\text{same distribution}
}
\]

in complete generality.

Conditions such as Carleman's criterion provide sufficient conditions
for moment determinacy.

%%%%%%%%%%%%%%%%%%%%%%%%%%%%%%%%%%%%%%%%%%%%%%%%%%%%%%%%%%%%
\subsection{Multivariate Moment Generating Functions}
\label{subsec:multivariate-mgf}
%%%%%%%%%%%%%%%%%%%%%%%%%%%%%%%%%%%%%%%%%%%%%%%%%%%%%%%%%%%%

For a random vector

\[
\mathbf{X}
=
(X_1,\ldots,X_n)^T,
\]

define

\[
\boxed{
M_{\mathbf{X}}(\mathbf{t})
=
\mathbb{E}
\left[
e^{\mathbf{t}^T\mathbf{X}}
\right].
}
\]

Here,

\[
\mathbf{t}
=
(t_1,\ldots,t_n)^T.
\]

Mixed partial derivatives recover joint moments.

%%%%%%%%%%%%%%%%%%%%%%%%%%%%%%%%%%%%%%%%%%%%%%%%%%%%%%%%%%%%
\subsection{Multivariate Characteristic Functions}
\label{subsec:multivariate-characteristic-function}
%%%%%%%%%%%%%%%%%%%%%%%%%%%%%%%%%%%%%%%%%%%%%%%%%%%%%%%%%%%%

For a random vector \(\mathbf{X}\),

\[
\boxed{
\varphi_{\mathbf{X}}(\mathbf{t})
=
\mathbb{E}
\left[
e^{i\mathbf{t}^T\mathbf{X}}
\right].
}
\]

This function uniquely determines the joint distribution of
\(\mathbf{X}\).

Thus characteristic functions extend naturally to multivariate
probability.

%%%%%%%%%%%%%%%%%%%%%%%%%%%%%%%%%%%%%%%%%%%%%%%%%%%%%%%%%%%%
\subsection{Independence and Multivariate Transforms}
\label{subsec:independence-multivariate-transforms}
%%%%%%%%%%%%%%%%%%%%%%%%%%%%%%%%%%%%%%%%%%%%%%%%%%%%%%%%%%%%

Random variables

\[
X_1,\ldots,X_n
\]

are independent if their joint characteristic function factors:

\[
\boxed{
\varphi_{X_1,\ldots,X_n}
(t_1,\ldots,t_n)
=
\prod_{j=1}^{n}
\varphi_{X_j}(t_j).
}
\]

This provides a transform-based characterization of independence.

%%%%%%%%%%%%%%%%%%%%%%%%%%%%%%%%%%%%%%%%%%%%%%%%%%%%%%%%%%%%
\subsection{Transform Strategy}
\label{subsec:probability-transform-strategy}
%%%%%%%%%%%%%%%%%%%%%%%%%%%%%%%%%%%%%%%%%%%%%%%%%%%%%%%%%%%%

When considering a probability transform, ask:

\begin{enumerate}
    \item Is the variable nonnegative integer-valued?
    \item Do I need coefficients or factorial moments?
    \item Does the MGF exist near zero?
    \item Would a characteristic function be safer?
    \item Is the variable nonnegative, making a Laplace transform natural?
    \item Am I adding independent random variables?
    \item Would a logarithm simplify multiplication into addition?
    \item Do I need moment information or the full distribution?
    \item Is an inversion formula necessary?
    \item Is the goal a limit theorem or asymptotic approximation?
\end{enumerate}

%%%%%%%%%%%%%%%%%%%%%%%%%%%%%%%%%%%%%%%%%%%%%%%%%%%%%%%%%%%%
\subsection{Common Errors}
\label{subsec:probability-transforms-common-errors}
%%%%%%%%%%%%%%%%%%%%%%%%%%%%%%%%%%%%%%%%%%%%%%%%%%%%%%%%%%%%

\subsubsection{Confusing PGF and MGF}

The PGF is

\[
\boxed{
G_X(s)
=
\mathbb{E}[s^X].
}
\]

The MGF is

\[
\boxed{
M_X(t)
=
\mathbb{E}[e^{tX}].
}
\]

They are different transforms.

%%%%%%%%%%%%%%%%%%%%%%%%%%%%%%%%%%%%%%%%%%%%%%%%%%%%%%%%%%%%
\subsubsection{Using PGFs for Arbitrary Real-Valued Variables}

PGFs are naturally suited to nonnegative integer-valued random
variables.

General real-valued variables are better handled by MGFs,
characteristic functions, or other transforms.

%%%%%%%%%%%%%%%%%%%%%%%%%%%%%%%%%%%%%%%%%%%%%%%%%%%%%%%%%%%%
\subsubsection{Assuming an MGF Always Exists}

MGFs may diverge for every nonzero \(t\) near zero.

Characteristic functions always exist.

%%%%%%%%%%%%%%%%%%%%%%%%%%%%%%%%%%%%%%%%%%%%%%%%%%%%%%%%%%%%
\subsubsection{Forgetting the Domain of an MGF}

For example, if

\[
X\sim\operatorname{Exp}(\lambda),
\]

then

\[
M_X(t)
=
\frac{\lambda}{\lambda-t}
\]

only for

\[
t<\lambda.
\]

The convergence domain matters.

%%%%%%%%%%%%%%%%%%%%%%%%%%%%%%%%%%%%%%%%%%%%%%%%%%%%%%%%%%%%
\subsubsection{Forgetting Independence When Multiplying Transforms}

The identity

\[
M_{X+Y}(t)
=
M_X(t)M_Y(t)
\]

requires independence.

The same warning applies to PGFs, Laplace transforms, and characteristic
functions of sums.

%%%%%%%%%%%%%%%%%%%%%%%%%%%%%%%%%%%%%%%%%%%%%%%%%%%%%%%%%%%%
\subsubsection{Confusing Raw Moments and Factorial Moments}

PGF derivatives naturally give factorial moments:

\[
G_X''(1)
=
\mathbb{E}[X(X-1)].
\]

MGF derivatives give raw moments:

\[
M_X''(0)
=
\mathbb{E}[X^2].
\]

%%%%%%%%%%%%%%%%%%%%%%%%%%%%%%%%%%%%%%%%%%%%%%%%%%%%%%%%%%%%
\subsubsection{Forgetting the Complex Unit in a Characteristic Function}

The characteristic function is

\[
\boxed{
\varphi_X(t)
=
\mathbb{E}[e^{itX}],
}
\]

not

\[
\mathbb{E}[e^{tX}].
\]

%%%%%%%%%%%%%%%%%%%%%%%%%%%%%%%%%%%%%%%%%%%%%%%%%%%%%%%%%%%%
\subsubsection{Assuming Characteristic Functions Are Real-Valued}

Characteristic functions may be complex-valued.

For symmetric distributions centered at zero, they are often real, but
this is not true in general.

%%%%%%%%%%%%%%%%%%%%%%%%%%%%%%%%%%%%%%%%%%%%%%%%%%%%%%%%%%%%
\subsubsection{Confusing Laplace Transform Sign Convention}

A probability Laplace transform is commonly written

\[
\boxed{
L_X(s)
=
\mathbb{E}[e^{-sX}].
}
\]

The negative sign distinguishes it from the MGF.

%%%%%%%%%%%%%%%%%%%%%%%%%%%%%%%%%%%%%%%%%%%%%%%%%%%%%%%%%%%%
\subsubsection{Differentiating Without Checking Moment Existence}

Formal derivatives of a transform correspond to moments only when the
required interchange of differentiation and expectation is justified.

%%%%%%%%%%%%%%%%%%%%%%%%%%%%%%%%%%%%%%%%%%%%%%%%%%%%%%%%%%%%
\subsubsection{Assuming Equal Moments Always Imply Equal Distributions}

Moment sequences do not uniquely determine every probability
distribution.

Characteristic functions, however, do uniquely determine probability
laws.

%%%%%%%%%%%%%%%%%%%%%%%%%%%%%%%%%%%%%%%%%%%%%%%%%%%%%%%%%%%%
\subsection{Applications}
\label{subsec:probability-transforms-applications}
%%%%%%%%%%%%%%%%%%%%%%%%%%%%%%%%%%%%%%%%%%%%%%%%%%%%%%%%%%%%

\begin{applicationbox}

\textbf{Distribution Sums.}
Transforms replace convolution by multiplication for independent sums.

\medskip

\textbf{Combinatorics.}
PGFs connect probability sequences with ordinary generating functions.

\medskip

\textbf{Statistics.}
MGFs and characteristic functions are used to derive sampling
distributions and asymptotic results.

\medskip

\textbf{Central Limit Theory.}
Characteristic functions provide one of the standard proof methods for
central limit theorems.

\medskip

\textbf{Queueing Theory.}
PGFs and Laplace transforms describe queue lengths, service times, and
waiting times.

\medskip

\textbf{Reliability Theory.}
Laplace transforms simplify sums of independent lifetime stages.

\medskip

\textbf{Stochastic Processes.}
Transform methods are fundamental for Poisson processes, branching
processes, and Markov models.

\medskip

\textbf{Signal Processing.}
Characteristic functions connect probability distributions with Fourier
analysis.

\medskip

\textbf{Compound Distributions.}
Transform composition provides a natural description of random sums and
compound Poisson models.

\end{applicationbox}

%%%%%%%%%%%%%%%%%%%%%%%%%%%%%%%%%%%%%%%%%%%%%%%%%%%%%%%%%%%%
\subsection{Exercises}
\label{subsec:probability-transforms-exercises}
%%%%%%%%%%%%%%%%%%%%%%%%%%%%%%%%%%%%%%%%%%%%%%%%%%%%%%%%%%%%

\begin{exercise}[1]
Define the probability generating function.
\end{exercise}

\begin{exercise}[1]
For a PGF \(G_X\), what is

\[
G_X(1)?
\]
\end{exercise}

\begin{exercise}[1]
Show that

\[
G_X(0)=P(X=0).
\]
\end{exercise}

\begin{exercise}[1]
Define the moment generating function.
\end{exercise}

\begin{exercise}[1]
Define the characteristic function.
\end{exercise}

\begin{exercise}[1]
Define the Laplace transform of a nonnegative random variable.
\end{exercise}

\begin{exercise}[1]
Define the cumulant generating function.
\end{exercise}

\begin{exercise}[1]
State the relationship between the PGF and MGF for a nonnegative
integer-valued random variable.
\end{exercise}

\begin{exercise}[1]
State the product rule for PGFs of independent sums.
\end{exercise}

\begin{exercise}[1]
State the product rule for MGFs of independent sums.
\end{exercise}

\begin{exercise}[2]
Find the PGF of

\[
X\sim\operatorname{Bernoulli}(p).
\]
\end{exercise}

\begin{exercise}[2]
Find the PGF of

\[
X\sim\operatorname{Bin}(n,p).
\]
\end{exercise}

\begin{exercise}[2]
Find the PGF of

\[
X\sim\operatorname{Poisson}(\lambda).
\]
\end{exercise}

\begin{exercise}[2]
Use the PGF of a Poisson variable to recover

\[
P(X=0).
\]
\end{exercise}

\begin{exercise}[2]
Use the PGF of a binomial variable to compute its mean.
\end{exercise}

\begin{exercise}[2]
Use PGF derivatives to derive the variance of a Poisson variable.
\end{exercise}

\begin{exercise}[2]
Find the MGF of a Bernoulli random variable.
\end{exercise}

\begin{exercise}[2]
Find the MGF of a binomial random variable.
\end{exercise}

\begin{exercise}[2]
Find the MGF of a Poisson random variable.
\end{exercise}

\begin{exercise}[2]
Find the MGF of

\[
X\sim\operatorname{Exp}(\lambda).
\]
\end{exercise}

\begin{exercise}[2]
State the convergence domain of the exponential MGF.
\end{exercise}

\begin{exercise}[2]
Find the characteristic function of a Bernoulli variable.
\end{exercise}

\begin{exercise}[2]
Find the characteristic function of a Poisson variable.
\end{exercise}

\begin{exercise}[2]
Find the Laplace transform of

\[
X\sim\operatorname{Exp}(\lambda).
\]
\end{exercise}

\begin{exercise}[3]
Show that

\[
G_X'(1)
=
\mathbb{E}[X].
\]
\end{exercise}

\begin{exercise}[3]
Show that

\[
G_X''(1)
=
\mathbb{E}[X(X-1)].
\]
\end{exercise}

\begin{exercise}[3]
Derive the variance formula

\[
\operatorname{Var}(X)
=
G_X''(1)
+
G_X'(1)
-
[G_X'(1)]^2.
\]
\end{exercise}

\begin{exercise}[3]
Use PGFs to show that the sum of independent binomial variables with
common success probability \(p\) is binomial.
\end{exercise}

\begin{exercise}[3]
Use PGFs to prove that the sum of independent Poisson random variables
is Poisson.
\end{exercise}

\begin{exercise}[3]
Use MGFs to derive the mean and variance of a normal random variable.
\end{exercise}

\begin{exercise}[3]
Use MGFs to prove that independent normal variables are closed under
addition.
\end{exercise}

\begin{exercise}[3]
Use MGFs to prove that independent chi-square random variables add their
degrees of freedom.
\end{exercise}

\begin{exercise}[3]
Use Laplace transforms to show that independent gamma variables with a
common rate add their shape parameters.
\end{exercise}

\begin{exercise}[3]
Show that

\[
|\varphi_X(t)|\leq1.
\]
\end{exercise}

\begin{exercise}[3]
Show that

\[
\varphi_X(-t)
=
\overline{\varphi_X(t)}.
\]
\end{exercise}

\begin{exercise}[3]
Derive the characteristic function of

\[
X\sim N(\mu,\sigma^2).
\]
\end{exercise}

\begin{exercise}[3]
Show that the characteristic function of a symmetric random variable
about zero is real-valued.
\end{exercise}

\begin{exercise}[3]
Prove that

\[
M_{aX+b}(t)
=
e^{bt}M_X(at).
\]
\end{exercise}

\begin{exercise}[3]
Prove that

\[
\varphi_{aX+b}(t)
=
e^{ibt}\varphi_X(at).
\]
\end{exercise}

\begin{exercise}[3]
For a Poisson variable, derive all cumulants from the cumulant generating
function.
\end{exercise}

\begin{exercise}[3]
Show that all cumulants of order at least \(3\) vanish for a normal
distribution.
\end{exercise}

\begin{exercise}[3]
Prove that cumulants add for independent sums.
\end{exercise}

\begin{exercise}[3]
Suppose \(N\) and the \(X_i\) are as in the random-sum construction.
Derive

\[
G_{S_N}(s)
=
G_N(G_X(s)).
\]
\end{exercise}

\begin{exercise}[3]
Derive the PGF of a compound Poisson random sum.
\end{exercise}

\begin{exercise}[4]
Study Lévy's continuity theorem and its use in proving convergence in
distribution.
\end{exercise}

\begin{exercise}[4]
Use characteristic functions to prove a version of the central limit
theorem for i.i.d. variables with finite variance.
\end{exercise}

\begin{exercise}[4]
Study Fourier inversion for probability densities.
\end{exercise}

\begin{exercise}[4]
Investigate inversion formulas for lattice distributions using
characteristic functions.
\end{exercise}

\begin{exercise}[4]
Study Carleman's condition for moment determinacy.
\end{exercise}

\begin{exercise}[4]
Investigate a moment-indeterminate distribution.
\end{exercise}

\begin{exercise}[4]
Study multivariate characteristic functions and prove their uniqueness.
\end{exercise}

\begin{exercise}[4]
Investigate joint cumulants and their connection with independence.
\end{exercise}

\begin{exercise}[4]
Study transform methods for branching processes.
\end{exercise}

\begin{exercise}[4]
Investigate Laplace--Stieltjes transforms in renewal and queueing theory.
\end{exercise}

%%%%%%%%%%%%%%%%%%%%%%%%%%%%%%%%%%%%%%%%%%%%%%%%%%%%%%%%%%%%
\subsection{Conceptual Summary}
\label{subsec:probability-transforms-summary}
%%%%%%%%%%%%%%%%%%%%%%%%%%%%%%%%%%%%%%%%%%%%%%%%%%%%%%%%%%%%

For a nonnegative integer-valued random variable,

\[
\boxed{
G_X(s)
=
\mathbb{E}[s^X]
}
\]

is the probability generating function.

Its coefficients are probabilities:

\[
\boxed{
P(X=k)
=
[s^k]G_X(s).
}
\]

Its derivatives give factorial moments:

\[
\boxed{
G_X'(1)=\mathbb{E}[X]
}
\]

and

\[
\boxed{
G_X''(1)=\mathbb{E}[X(X-1)].
}
\]

The moment generating function is

\[
\boxed{
M_X(t)
=
\mathbb{E}[e^{tX}].
}
\]

When it exists near zero,

\[
\boxed{
M_X^{(k)}(0)
=
\mathbb{E}[X^k].
}
\]

The characteristic function is

\[
\boxed{
\varphi_X(t)
=
\mathbb{E}[e^{itX}].
}
\]

It always exists and uniquely determines the distribution.

For independent variables,

\[
\boxed{
\varphi_{X+Y}(t)
=
\varphi_X(t)\varphi_Y(t).
}
\]

For nonnegative variables, the Laplace transform is

\[
\boxed{
L_X(s)
=
\mathbb{E}[e^{-sX}].
}
\]

The cumulant generating function is

\[
\boxed{
K_X(t)
=
\ln M_X(t).
}
\]

Its derivatives produce cumulants:

\[
\boxed{
\kappa_1=\mathbb{E}[X],
\qquad
\kappa_2=\operatorname{Var}(X).
}
\]

For independent sums,

\[
\boxed{
\kappa_k(X+Y)
=
\kappa_k(X)+\kappa_k(Y).
}
\]

Transforms therefore convert

\[
\boxed{
\text{convolution}
\longrightarrow
\text{multiplication},
}
\]

and encode

\[
\boxed{
\text{probabilities}
+
\text{moments}
+
\text{distributional structure}
+
\text{asymptotic behavior}.
}
\]

%%%%%%%%%%%%%%%%%%%%%%%%%%%%%%%%%%%%%%%%%%%%%%%%%%%%%%%%%%%%
\subsection{Section Connections}
\label{subsec:probability-transforms-connections}
%%%%%%%%%%%%%%%%%%%%%%%%%%%%%%%%%%%%%%%%%%%%%%%%%%%%%%%%%%%%

\chapterconnection{
Probability Transforms provides an algebraic and analytic counterpart to
the direct transformation and convolution methods of
Section~\ref{sec:functions-random-variables}. Probability generating
functions connect directly with the generating functions of Chapter 6,
while characteristic functions connect probability with Fourier
analysis. Moment and cumulant generating functions encode the moments
developed in Section~\ref{sec:expectation-moments}. The next section,
Laws of Large Numbers, studies convergence of random averages, while
Central Limit Theorems use characteristic functions and small-argument
transform expansions to analyze standardized sums. Later sections on
stochastic processes use PGFs, Laplace transforms, and characteristic
functions extensively in Poisson, branching, renewal, queueing, and
diffusion models.
}

%%%%%%%%%%%%%%%%%%%%%%%%%%%%%%%%%%%%%%%%%%%%%%%%%%%%%%%%%%%%
% SECTION SOURCE TRACKING
%%%%%%%%%%%%%%%%%%%%%%%%%%%%%%%%%%%%%%%%%%%%%%%%%%%%%%%%%%%%

%------------------------------------------------------------
% 7.8 Probability Transforms
%------------------------------------------------------------
%
% Source basis:
%   - Standard probability generating-function theory.
%   - Standard moment-generating-function theory.
%   - Standard characteristic-function theory.
%   - Standard Laplace-transform probability methods.
%   - Standard cumulant theory.
%
% Used for:
%   - probability generating functions;
%   - coefficient extraction;
%   - PGF normalization;
%   - Bernoulli PGFs;
%   - binomial PGFs;
%   - geometric PGFs;
%   - Poisson PGFs;
%   - factorial moments;
%   - variance from PGFs;
%   - independent sums using PGFs;
%   - moment generating functions;
%   - moments from MGFs;
%   - Bernoulli MGFs;
%   - binomial MGFs;
%   - Poisson MGFs;
%   - normal MGFs;
%   - exponential MGFs;
%   - gamma MGFs;
%   - chi-square MGFs;
%   - independent sums using MGFs;
%   - MGF uniqueness;
%   - failure of MGF existence;
%   - characteristic functions;
%   - Euler representation;
%   - existence and boundedness;
%   - Bernoulli characteristic functions;
%   - binomial characteristic functions;
%   - Poisson characteristic functions;
%   - normal characteristic functions;
%   - Cauchy characteristic functions;
%   - characteristic functions of sums;
%   - uniqueness;
%   - Fourier-transform interpretation;
%   - inversion preview;
%   - Lévy continuity theorem preview;
%   - Laplace transforms;
%   - exponential and gamma Laplace transforms;
%   - cumulant generating functions;
%   - cumulants;
%   - cumulant additivity;
%   - Poisson and normal cumulants;
%   - transform relationships;
%   - random sums;
%   - compound Poisson preview;
%   - branching-process preview;
%   - queueing and reliability connections;
%   - transform methods for limit theorems;
%   - moment determinacy preview;
%   - multivariate transforms;
%   - applications and exercises.
%
% Notes:
%   - Laws of Large Numbers are developed in Section 7.9.
%   - Central Limit Theorems are developed in Section 7.10.
%   - Fourier analysis is developed independently in Chapter 5.
%   - Measure-theoretic probability gives the rigorous framework for
%     characteristic functions, convergence, and transform inversion.
%   - Stochastic processes later make extensive use of PGFs and Laplace
%     transforms.
%
%%%%%%%%%%%%%%%%%%%%%%%%%%%%%%%%%%%%%%%%%%%%%%%%%%%%%%%%%%%%